% Standalone source-record documentation; no research attempts.
\documentclass[11pt]{article}
\usepackage[margin=25mm]{geometry}
\usepackage{fontspec,unicode-math}
\setmainfont{Latin Modern Roman}
\setmathfont{Latin Modern Math}
\usepackage{amsmath,amssymb,mathtools,xurl,hyperref}
\hypersetup{colorlinks=true,urlcolor=blue,linkcolor=blue}
\setcounter{secnumdepth}{0}
\setlength{\parindent}{0pt}
\setlength{\parskip}{5pt}
\setlength{\emergencystretch}{3em}
\newcommand{\sref}[2]{\hyperlink{source:#1:#2}{[S#2]}}
\newcommand{\cataloguescope}[1]{\paragraph{Source record.}#1}
\newcommand{\floor}[1]{\left\lfloor#1\right\rfloor}
\begin{document}
% UnsolvedMath ID 20001323; source catalogue code AIM-CRYPTOGRAPHY-0010
\setcounter{section}{20001322}
\section{A fiber--leakage--arity audit for group-theoretic cryptographic building blocks}\label{20001323}
{\small\sffamily UnsolvedMath ID 20001323 · Computer Science}
\subsection{Definitions and mathematical statement}

\paragraph{Shared notation.}Standard number systems are $\mathbb N,\mathbb Z,\mathbb Q,\mathbb R,\mathbb C$. Specialized notation and index ranges follow the cited source and the definitions below.


A one-way function is efficiently computable while inversion on its specified input distribution is computationally hard. A multilinear map is linear or homomorphic separately in each argument in its stated algebraic domains. A cryptographic construction from a group problem needs an explicit representation, hardness assumption, and reduction connecting that assumption to the desired security property.
\paragraph{Record scope and qualifications.}
The catalogue title is retained as metadata. Named certificates, contracts, bounds or frameworks in that title are not asserted as results of the original source extract.
No underlying problem, map, distribution, or security reduction is given. The record states a construction agenda without asserting that arbitrary group-theoretic hardness yields those primitives. The supplied catalogue marks this record partially solved; its present research status was not independently checked.
The source text below is rendered literally where necessary; damaged formulas and composite extracts are not silently repaired into a new mathematical claim.
\paragraph{Statement recorded in the supplied source.}
Cryptographic building blocks from problems in group theory

We would like to use computational problems in group theory to derive fundamental cryptographic objects including one-way functions and multilinear maps. \sref{20001323}{2}
\paragraph{Source status.}
The supplied catalogue labels this source record partially solved; that metadata was checked against the supplied snapshot on 2026-10-11. The label is preserved for provenance and is not an independent verification of present open status or mathematical correctness.
\subsection{Short English statement}
The source asks for cryptographic building blocks derived from computational group problems.
\subsection{Sources}
\begin{itemize}
\item[\hypertarget{source:20001323:1}{S1}] Ulam AI and the UnsolvedMath Contributors, UnsolvedMath: A Curated Collection of Open Mathematics Problems, record 20001323 (AIM-CRYPTOGRAPHY-0010). Catalogue curation and metadata: CC BY 4.0. The source extract is retained; definitions, English summaries and qualifications are editorial additions.
\url{https://huggingface.co/datasets/ulamai/UnsolvedMath}.
\item[\hypertarget{source:20001323:2}{S2}] AIM Workshop Problem Lists, source linked by catalogue record AIM-CRYPTOGRAPHY-0010.
\url{http://aimpl.org/postquantgroup/3/}.
\end{itemize}
\paragraph{Preparation and review.}
The original source and UnsolvedMath contributors retain attribution for the source extract. Codex (GPT-6; exact serving revision not exposed) assisted with definitions, English summaries, scope notes and formatting. No new human mathematical authorship or independently verified solution is asserted. Permanent source ID 20001323 and code AIM-CRYPTOGRAPHY-0010 checked against the complete supplied registry and current repository. Distinct upstream records, including repeated questions and broad or damaged extracts, are retained for complete coverage of the source CS category. A repeated record is not a new independent conjecture. Complete mathematical-equivalence and current-literature review has not been performed. Preparation and final review follow contribution prompts C and D, with source gaps disclosed. Structural and rendering checks do not establish a new theorem or open status.
\label{end:20001323}
\end{document}
